\chapter{Measurement Design}
\label{ch:measurement}

The first study asked whether an extra coding check improved a prediction of success on a separate task. The language model
\texttt{gpt-oss-120b} supplied answers and code for 24 authored Python cases covering four concepts and eight misconceptions.
Each case began with an answer that reviewers could read. A related coding task supplied additional evidence. The system
also formed controls by substituting unrelated evidence or reversing test results. A separate coding task, called the criterion
task, supplied the outcome used to assess the predictions. The workflow saved predictions before requesting this final code
response. Its ordering checks are described in Chapter~\ref{ch:system}.

\section{Case Construction}

The corpus covers assignment evaluation, object aliasing, function arguments and conditional control. Each concept contains
two authored misconception groups, with three cases per group. Examples include treating assignment as copying a mutable
object, assuming that a function copies its argument, and treating a conditional as two independent checks. These groups
define the intended coverage. Their frequency in the corpus does not estimate prevalence among students.

Each case links a public prompt to an executable evidence task and a separate criterion task. The manifest records their
paths and hashes, an initial tracker state, control assignments, and evaluator-only review and structural-difference records.
The structural records explain how the evidence and criterion tasks differ. For example, one evidence task asks for the sorted
result after adding a permission, while its criterion asks for a removal operation through a function over a different
collection. Differences in operation, signature and values reduce literal answer reuse. They do not by themselves validate
educational transfer or exclude shared shortcuts.

The manifest also retains authored evidence-pattern labels used during benchmark construction. These are design metadata,
not observed performance labels for the external model. Its public response is rated independently, and its criterion outcome
comes from executing its submitted code. Neither outcome is inferred from the intended case pattern. Calling the corpus
held-out refers to its experimental role. It does not establish that these concepts or similar questions were absent from
the model's training data.

\section{Control Construction}

The unrelated control uses a preassigned source case whose target-concept label differs from the target case. It copies that
source's recorded evidence pass and fail counts unchanged. The implementation rejects self-assignment and same-concept
assignment, and records the source summary and assignment hashes. The mapping is fixed rather than selected after inspecting
which source gives a favourable comparison. Different concept labels define unrelated evidence for this experiment.
The underlying programming skills may still overlap.

The corrupted control swaps the target evidence's pass and fail counts. All-pass evidence becomes all-fail evidence, and
vice versa. A mixed result remains in the conflicting category after the swap, so corruption does not necessarily change
the tracker's update for every possible input. The recorded seed is provenance metadata. This transformation is deterministic
and does not draw random flips. It changes the summary supplied to the tracker, not the submitted program or the sandbox.

Both controls reuse existing evidence summaries. They require no extra model generation and are not separate tutor
conversations. This design isolates the response of the fixed tracker to substituted evidence. It cannot show how a live
tutor or student would react to receiving an unrelated question or misleading feedback.

\section{From Observations to Predictions}

The evaluation model answered the public, evidence and criterion prompts in separate contexts. Here, isolated means that a
response did not receive the conversation history from another channel. It does not mean that the model had never encountered
similar tasks during training. The later criterion was withheld from the tracker and prediction process. Its completion tests
model performance rather than learning caused by the earlier prompt.

The v1 tracker is an authored additive rule, separate from the Bayesian tracker introduced later. It starts at a score of 0.5.
A correct observation adds 0.2. An incorrect answer or misconception subtracts 0.2. A conflicting observation subtracts 0.1,
while an uncertain or empty observation leaves the score unchanged. Each update is restricted to the interval $[0,1]$:
\begin{equation}
  s^{+}=\min(1,\max(0,s+\Delta_c)),
\end{equation}
where $\Delta_c$ is the fixed change for the observation category. The stored field is named
\texttt{mastery\_probability}, but the score is not a calibrated probability or a measurement of human mastery. Recorded
confidence values do not affect this update. The binary prediction is success when the final score is at least 0.6.

Independent public-answer ratings provide the first observation. Each condition then starts from an identical copy of the
post-answer state. Dialogue alone receives no further update. The other conditions receive the category derived from their
respective test counts: all checks passing means correct, all failing means incorrect, and a mixture means conflicting.
An empty set of test results is an invalid input. The conditions are evaluated separately. Updates from
one condition cannot carry into another.

This rule deliberately compresses the evidence. It does not inspect program meaning or weight a result by its relevance to
the target concept. Two checks with the same outcome category therefore produce the same update, even when one is unrelated.
The unrelated control tests this limitation of the chosen prediction rule. Equality between those conditions cannot establish
that relevance never matters for a richer tracker.

\section{Public-Answer Review}

Two reviewers assessed the 24 public responses: an MSc Computing student and a senior software engineer. Each received a CSV
containing the case identifier, public prompt, visible response and response hash, together with the shared category guide.
The packet excluded executable probes, test results, criterion material, authored evidence patterns and condition predictions.
